\section*{Capítulo \ref{ch:g11:drift-to-spreading} --- De la deriva continental a la expansión del fondo oceánico}

\begin{solution}{exo:g11:drift-to-spreading:1}
La \omterm{def:g11:drift-to-spreading:spreading}{expansión del fondo oceánico} es la fabricación de fondo oceánico nuevo en
la cresta de una dorsal oceánica, donde la roca fundida sube y se solidifica,
seguida de su alejamiento de la dorsal hacia ambos lados. La velocidad total
es la velocidad a la que se separan las dos placas; la semivelocidad es la
velocidad a la que el fondo de un lado se aleja del eje. Si la expansión es
simétrica, velocidad total $= 2
\times$ semivelocidad.
\end{solution}

\begin{solution}{exo:g11:drift-to-spreading:2}
Por encima de su \omterm{def:g11:drift-to-spreading:remanence}{temperatura de Curie}, unos \qty{580}{\celsius}, la
magnetita no puede conservar una magnetización propia; solo al enfriarse por
debajo de ella sus granos se magnetizan según el campo y conservan esa
dirección. Si se vuelve a calentar a \qty{700}{\celsius}, por encima de la
\omterm{def:g11:drift-to-spreading:remanence}{temperatura de Curie}, la roca pierde su antiguo registro; al enfriarse de
nuevo, registra el campo de esa nueva época.
\end{solution}

\begin{solution}{exo:g11:drift-to-spreading:3}
El sonido baja y vuelve: $d = 1500 \times 5.2/2 = \qty{3900}{m}$.
\end{solution}

\begin{solution}{exo:g11:drift-to-spreading:4}
El intervalo normal actual empezó hace \num{0.78} millones de años: ha durado
\num{780000} años. El episodio Jaramillo duró $1.07 - 0.99 = \num{0.08}$
millones de años, \num{80000} años. Contando los límites de menos de cinco
millones de años (0.78, 0.99, 1.07, 1.77, 1.95, 2.14, 2.15, 2.581, 3.04, 3.11,
3.22, 3.33, 3.58, 4.18, 4.29, 4.48, 4.62, 4.80, 4.89, 4.98): veinte inversiones.
\end{solution}

\begin{solution}{exo:g11:drift-to-spreading:5}
$65\,000/40\,000 \approx 1.6$: las dorsales darían alrededor de vuelta y
media a la Tierra.
\end{solution}

\begin{solution}{exo:g11:drift-to-spreading:6}
$20/0.78 \approx \qty{25.6}{km}$ por millón de años, es decir, una
semivelocidad de unos \qty{2.6}{cm} al año; la velocidad total es de unos
\qty{5.1}{cm} al año.
\end{solution}

\begin{solution}{exo:g11:drift-to-spreading:7}
Oeste: $45/3.58 \approx 12.6$~km por millón de años, \qty{1.26}{cm} al año;
este: $44/3.58 \approx 12.3$~km por millón de años, \qty{1.23}{cm} al año.
Velocidad total: $1.26 + 1.23 \approx \qty{2.5}{cm}$ al año. Las dos
semivelocidades difieren en un \qty{2}{\percent} más o menos: la expansión es
casi exactamente simétrica.
\end{solution}

\begin{solution}{exo:g11:drift-to-spreading:8}
\qty{1.9}{cm} al año son \qty{19}{km} por millón de años: $990/19 \approx 52$
millones de años, algo más que los 49 millones de años del sedimento más
antiguo (alrededor de un \qty{6}{\percent}).
\end{solution}

\begin{solution}{exo:g11:drift-to-spreading:9}
A \qty{20}{km} por millón de años: \qty{15.6}{km}, \qty{19.8}{km},
\qty{21.4}{km}, \qty{35.4}{km} y \qty{39.0}{km}. Desde el eje: normal de 0
a \qty{15.6}{km}, inversa hasta \qty{19.8}{km}, normal hasta \qty{21.4}{km}
(Jaramillo), inversa hasta \qty{35.4}{km}, normal hasta \qty{39.0}{km}
(Olduvai), y después inversa.
\end{solution}

\begin{solution}{exo:g11:drift-to-spreading:10}
$64 \times 10^6/4.6 \times 10^9 \approx \num{0.014}$, alrededor del
\qty{1.4}{\percent} de la edad de la Tierra. Apoya el argumento de Hess según
el cual el fondo oceánico es joven porque se fabrica y se destruye sin cesar:
ningún fondo oceánico ha sobrevivido a la mayor parte de la historia de la
Tierra.
\end{solution}

\begin{solution}{exo:g11:drift-to-spreading:11}
En cada momento, el fondo se forma en el eje con la polaridad de ese momento,
y se divide: una mitad se aleja hacia cada lado a la misma velocidad. Por
tanto, dos franjas formadas al mismo tiempo están a la misma distancia del
eje, una a cada lado, con la misma polaridad: cada lado es una copia del
otro, reflejada en el eje. La simetría se rompe si el fondo se aleja más
deprisa por un lado que por el otro (expansión asimétrica), o si el eje salta
a una nueva posición.
\end{solution}

\begin{solution}{exo:g11:drift-to-spreading:12}
(1) \qty{5}{cm} son $5 \times 10^{-5}$~\unit{km}: cada año se añade una franja
de $65\,000 \times 5 \times 10^{-5} = \qty{3.25}{km^2}$.
(2) $2.9 \times 10^8/3.25 \approx 9 \times 10^7$ años, unos 90 millones de
años. (3) Es del mismo orden que la edad del fondo más antiguo (190 millones
de años) y más corto: para que el fondo oceánico siga siendo tan joven, en
algún lugar tiene que destruirse tanto fondo como se fabrica en las dorsales.
\end{solution}

\begin{solution}{exo:g11:drift-to-spreading:13}
El episodio duró \num{0.08} millones de años. A \qty{12.5}{km} por millón de
años, su banda mide $0.08 \times 12.5 = \qty{1}{km}$ de ancho; a \qty{75}{km}
por millón de años, $0.08 \times 75 = \qty{6}{km}$. Vista desde la superficie
del mar, \qty{3}{km} más arriba, una banda mucho más estrecha que su
profundidad solo produce un campo débil, que se confunde con los campos de
sus anchas vecinas de polaridad contraria; una banda el doble de ancha que su
profundidad destaca.
\end{solution}

\begin{solution}{exo:g11:drift-to-spreading:14}
Punto 16: $176/11 \approx \qty{16}{m}$ por millón de años; punto 20: $72/67
\approx \qty{1.1}{m}$ por millón de años. El sedimento se acumula a ritmos muy
distintos de un lugar a otro y de una época a otra, así que su espesor no da
la edad del fondo. La edad la dan los fósiles del sedimento más antiguo, que
descansa directamente sobre el basalto.
\end{solution}

\begin{solution}{exo:g11:drift-to-spreading:15}
(1) Oeste: $58/2.581 \approx \qty{22.5}{km}$ por millón de años,
\qty{2.25}{cm} al año; este: $70/2.581 \approx \qty{27.1}{km}$ por millón de
años, \qty{2.71}{cm} al año; velocidad total: $128/2.581 \approx \qty{49.6}{km}$ por
millón de años, unos \qty{5.0}{cm} al año. (2) Unos \qty{225}{km} al oeste y
\qty{271}{km} al este. (3) Los dos límites estaban juntos en el eje cuando se
produjo la inversión; desde entonces se han separado a la velocidad total,
se reparta como se reparta el movimiento entre los dos lados. Su separación
dividida entre su edad es la velocidad total.
\end{solution}

\begin{solution}{pb:g11:drift-to-spreading:1}
\textbf{1.} El punto 16 es el más cercano (\qty{191}{km}, 11 millones de
años), y el punto 20 el más lejano (\qty{1270}{km}, 67 millones de años):
cuanto más lejos de la dorsal, más antiguo es el fondo.

\textbf{2.} El sedimento empieza a depositarse sobre el basalto en cuanto el
basalto se ha formado, así que la capa más antigua tiene más o menos la edad
del basalto. Solo puede ser algo más joven, nunca más antigua.

\textbf{3.} En kilómetros por millón de años: punto 16: $191/11 \approx 17.4$;
punto 15: $380/24 \approx 15.8$; punto 18: $506/26 \approx 19.5$; punto 17:
$643/33 \approx 19.5$; punto 14: $727/40 \approx 18.2$; punto 19: $990/49
\approx 20.2$; punto 20: $1270/67 \approx 19.0$.

\textbf{4.} $129.5/7 \approx \num{18.5}$~km por millón de años, unos
\qty{1.85}{cm} al año.

\textbf{5.} Sí: sus cocientes (19.5) están entre los de los puntos del oeste.
El fondo se ha expandido a más o menos la misma velocidad por los dos lados:
la expansión es simétrica.

\textbf{6.} $t = d/u$: la edad es proporcional a la distancia.

\textbf{7.} $u = 208\,276/10\,952 \approx \num{19.0}$~km por millón de años,
es decir, \qty{1.90}{cm} al año.

\textbf{8.} Unos $2 \times 1.9 = \qty{3.8}{cm}$ al año.

\textbf{9.} Sí: \qty{1.9}{cm} al año es algo menos de \qty{2}{cm} al año.

\textbf{10.} $1270/19.0 \approx 67$ millones de años, lo mismo que su
sedimento más antiguo.

\textbf{11.} $191/9 \approx \num{21.2}$~km por millón de años,
\qty{2.1}{cm} al año.

\textbf{12.} $1270/71 \approx \num{17.9}$~km por millón de años,
\qty{1.8}{cm} al año.

\textbf{13.} $180\,329/9854 \approx \num{18.3}$~km por millón de años,
\qty{1.83}{cm} al año.

\textbf{14.} $(1.90 - 1.83)/1.83 \approx \qty{4}{\percent}$: las dos
velocidades concuerdan dentro de las incertidumbres de las edades y de las
distancias.

\textbf{15.} Una edad procede de los fósiles del sedimento; la otra, del
magnetismo del basalto y de la escala de tiempo de las inversiones. Un error
de un método no aparecería en el otro: su concordancia es una prueba sólida
de la \omterm{def:g11:drift-to-spreading:spreading}{expansión del fondo oceánico}.

\textbf{16.} $13.8 + 13.8 = \ang{27.6}$.

\textbf{17.} $111.2 \times \cos 30^\circ \approx \qty{96.3}{km}$ por grado;
$27.6 \times 96.3 \approx \qty{2660}{km}$.

\textbf{18.} $2660/19.0 \approx 140$ millones de años.

\textbf{19.} $2660/130 \approx \num{20.5}$~km por millón de años, unos
\qty{2.0}{cm} al año, algo más que \qty{1.9}{cm} al año. Como los últimos 67
millones de años transcurrieron a unos \qty{1.9}{cm} al año, la apertura
anterior fue algo más rápida: el fondo entre el punto 20 y el borde,
$2660 - 1270 =
\qty{1390}{km}$, se formó en $130 - 67 = 63$ millones de años,
a unos \qty{2.2}{cm} al año.

\textbf{20.} \textbf{El fondo oceánico más antiguo junto a África a \ang{30} S
tiene unos 130 millones de años} (140 millones de años si se extrapola la
velocidad actual): el Atlántico Sur empezó a abrirse allí en el periodo
Cretácico.
\end{solution}
